\documentclass[11pt]{article}
\usepackage[letterpaper,margin=1in]{geometry}
\usepackage{amsmath,amssymb,amsthm,booktabs,microtype,xurl}
\usepackage[hidelinks]{hyperref}

\newtheorem{theorem}{Theorem}
\newtheorem{proposition}[theorem]{Proposition}
\newtheorem{lemma}[theorem]{Lemma}
\newtheorem{corollary}[theorem]{Corollary}
\newtheorem{computationallemma}[theorem]{Computational Lemma}
\newtheorem{computationalclaim}[theorem]{Computational Claim}
\newcommand{\Z}{\mathbb Z}
\newcommand{\PAF}{\operatorname{PAF}}
\newcommand{\BS}{\operatorname{BS}}
\newcommand{\TU}{\operatorname{TU}}
\newcommand{\dist}{\operatorname{dist}}

\title{Exact local obstructions around a\\
64-modular Hadamard matrix of order 668}
\author{Alec Kriebel\\[2pt]\large with heavy assistance from ChatGPT 5.6 Sol}
\date{Research draft prepared 22 July 2026\\Not peer reviewed}

\begin{document}
\maketitle

\begin{abstract}
Eliahou recently constructed a 64-modular Hadamard matrix of order $668$
from a structured length-$167$ Golay quadruple modulo $64$.  We study two
sharply delimited exact-repair questions around that construction.  First,
if Eliahou's sequence $q=(+^{83},-^2,+^{81},-)$ is held fixed, then no
choice of the other binary sequence makes the special quadruple exactly
Golay.  The proof is a parity telescope followed by even--odd decimation:
any such repair would produce a Turyn sequence in $\TU(41)$, which is empty.
We independently reproduce that endpoint by an exact 461-shard enumeration.
Second, allowing both sequences to vary is equivalent to the base-sequence
problem $\BS(84,83)$.  A finite decomposition by ordinary and alternating
margins, endpoint-quad parity, and primitive-root norm identities reports no
exact base sequence within raw labelled Hamming distance $18$ of the
published seed.  The latter is a local computational claim about one
orientation of one structured construction, not a nonexistence result for
$\BS(84,83)$ or for Hadamard matrices of order $668$.  Exact data,
bounded-memory generators, and independent artifact checkers are being
assembled for release with the paper; proof-grade certification of all
solver-infeasible leaves is a release requirement.
\end{abstract}

\noindent\textbf{Verification disclaimer.}
The human author is a complete amateur exploring the limits of AI-assisted
mathematics and cannot independently validate these claims.  This is an
unreviewed artifact for expert checking, not an established result.  Passing
the included programs establishes properties of their encoded finite models;
it is not peer review and does not by itself establish priority.

\section{Introduction}

A Hadamard matrix of order $n$ is a matrix
$H\in\{\pm1\}^{n\times n}$ satisfying
\[
 HH^\top=nI_n.
\]
The Hadamard conjecture asserts existence whenever $4\mid n$.  The smallest
order for which no matrix is presently known is $668$; the next unresolved
orders in the implemented range are $716$, $892$, and
$1132$~\cite{cati-pasechnik,sage}.

Eliahou constructed a $64$-modular matrix at order $668$, meaning
\[
 HH^\top\equiv668I_{668}\pmod {64},
\]
from four binary sequences of length $167$ and the Goethals--Seidel
array~\cite{eliahou-2025}.  A subsequent paper embeds this example in
families of 64-modular constructions~\cite{eliahou-2026}.  The order-$668$
seed is unusually close to exact complementarity: only thirteen positive
aperiodic correlation lags are nonzero.

The natural question is whether a small sign repair can turn this structured
seed into an exact Golay quadruple and hence into a true Hadamard matrix.  The
present paper proves one bounded negative result and reports one larger
solver-backed exclusion.

\begin{theorem}[Fixed-$q$ obstruction]\label{thm:fixedq}
Keep Eliahou's length-$167$ sequence
\[
 q=(+^{83},-^2,+^{81},-)
\]
and the special form
\[
 (s,s',sq,(sq)').
\]
No binary sequence $s$ makes this quadruple exactly Golay.
\end{theorem}

\begin{computationalclaim}[Raw radius-$18$ solver report]\label{thm:radius}
Let $\mathcal E\in\{\pm1\}^{334}$ be the labelled base quadruple obtained
from Eliahou's published order-$668$ seed.  The recorded CP-SAT decomposition
reports that there is no
\[
 (A;B;C;D)\in\BS(84,83)
\]
with raw labelled Hamming distance
\[
 \dist\bigl((A;B;C;D),\mathcal E\bigr)\le18.
\]
\end{computationalclaim}

The word \emph{raw} is essential.  No distance-changing equivalence is
applied, and no global neighbourhood of an equivalence class is claimed.
The ball has radius $18$ in a $334$-dimensional binary cube; even a certified
exclusion would not suggest that the unrestricted base-sequence problem is
nearly exhausted.

Theorem~\ref{thm:fixedq} is a short mathematical reduction whose endpoint
uses the exhaustive classification of Turyn sequences of long length below
$43$~\cite{edmondson}.  Claim~\ref{thm:radius} is a finite solver report.
Its decomposition and decoded root witnesses have dependency-free
checkers.  The public release is not proof-grade until every solver-reported
infeasible leaf is also covered by an independently replayable proof or a
second exact enumeration; Section~\ref{sec:certification} makes that
distinction explicit.

\section{Sequences, the modular seed, and base sequences}

For a sign sequence $X=(x_0,\ldots,x_{r-1})$, write
\[
 N_X(k)=\sum_{i=0}^{r-k-1}x_ix_{i+k}
 \qquad(0\le k<r)
\]
for its aperiodic autocorrelation.  A quadruple
$(X_1,X_2,X_3,X_4)$ of equal length is Golay when
\[
 \sum_{j=1}^4N_{X_j}(k)=0
 \qquad(1\le k<r).
\]

Let $h\in\{\pm1\}^{167}$ be $+1$ on coordinates $0,\ldots,83$ and
$-1$ on coordinates $84,\ldots,166$, and write $s'=hs$ for pointwise
multiplication.  Eliahou's quadruple is
\begin{equation}
 \mathcal G(s,q)=(s,hs,sq,hsq).
 \label{eq:special}
\end{equation}
The published $q$ has alternating run lengths
\[
 (83,2,81,1).
\]
For reproducibility, the published $s$ has run-length word
\begin{align*}
 &((4)^5,(2,1,1)^5,1,5,(4)^4,(2,1,1)^6,(4)^4,3,\\
 &\hspace{32mm}(1,2,1)^5,3,(4)^4,3,(1,2,1)^5),
\end{align*}
where an exponent means repetition of a block and the first run has sign
$+1$.  These two words are encoded directly in run-length form in the
accompanying checker.

The nonzero summed aperiodic correlations of the quadruple occur at the
thirteen lags shown below.
\[
\begin{array}{c|rrrrrrrrrrrrr}
k&4&8&12&16&26&30&34&38&42&46&50&54&58\\ \hline
\sum_jN_{X_j}(k)
&-512&384&-256&128&-64&128&-192&256&-320&256&-192&128&-64
\end{array}
\]
Thus every residual is divisible by $64$, but the quadruple is not exact.

A base sequence in $\BS(m,n)$ is a quadruple
$(A;B;C;D)$ with $|A|=|B|=m$, $|C|=|D|=n$, and
\[
 N_A(k)+N_B(k)+N_C(k)+N_D(k)=0
 \qquad(1\le k<m),
\]
where the shorter terms vanish once $k\ge n$.

\begin{proposition}[Exact change of variables]\label{prop:bs}
The special quadruple $\mathcal G(s,q)$ of length $167$ is Golay if and
only if
\[
 A=s_{[0,84)},\quad B=(sq)_{[0,84)},\quad
 C=s_{[84,167)},\quad D=(sq)_{[84,167)}
\]
belong to $\BS(84,83)$.  The correspondence is bijective, with inverse
\[
 s=A\Vert C,\qquad q=(AB)\Vert(CD).
\]
\end{proposition}

\begin{proof}
At lag $k$, the coefficient in the summed autocorrelation of
\eqref{eq:special} is
\[
 \sum_i s_is_{i+k}
 (1+h_ih_{i+k})(1+q_iq_{i+k}).
\]
Cross-half terms vanish.  The remaining expression is twice
\[
 N_A(k)+N_B(k)+N_C(k)+N_D(k).
\]
It vanishes automatically for $k\ge84$, proving the equivalence.  The
displayed inverse follows coordinatewise.
\end{proof}

Proposition~\ref{prop:bs} is useful in both directions: an exact repair gives
a standard base sequence, and any base-sequence candidate gives explicit
$(s,q)$ and a full order-$668$ Goethals--Seidel verification path.

\section{The fixed-\texorpdfstring{$q$}{q} obstruction}

We prove Theorem~\ref{thm:fixedq}.  The special shape of the four runs in
$q$ makes almost all variables disappear from the correlation equations.
Write
\[
 X=s_{[0,83)},\quad u=s_{84},\quad
 Y=s_{[85,166)},\quad v=s_{166}.
\]
The coordinate $s_{83}$ is isolated.  Direct inspection of the factor
$(1+h_ih_{i+k})(1+q_iq_{i+k})$ in
\eqref{eq:special} leaves the within-$X$ and within-$Y$ pairs at lags
$1\le k\le80$, only the two within-$X$ pairs at lag $81$, and the pair
$x_0x_{82}$ together with $uv$ at lag $82$.  No active pair remains at a
larger lag.  Exact complementarity is therefore equivalent to
\begin{align}
 N_X(k)+N_Y(k)&=0 &&(1\le k\le80),\label{eq:fixed-main}\\
 x_0x_{81}+x_1x_{82}&=0,\label{eq:fixed-81}\\
 x_0x_{82}+uv&=0.\label{eq:fixed-end}
\end{align}

\begin{lemma}[Parity telescope]\label{lem:telescope}
Every solution of \eqref{eq:fixed-main}--\eqref{eq:fixed-end} satisfies
\begin{align}
 x_{82-i}&=(-1)^{i+1}x_i &&(0\le i\le82),\label{eq:xreverse}\\
 y_{80-i}&=(-1)^iy_i &&(0\le i\le80),\label{eq:yreverse}\\
 u&=v.\label{eq:uv}
\end{align}
\end{lemma}

\begin{proof}
For $0\le k\le81$, let $P_k$ be the product of all sign terms in
$N_X(k)+N_Y(k)$, with the $Y$ product empty at $k=81$.  Every lag-zero term
is a square, so $P_0=1$.  At positive lag the zero sum contains
$2(82-k)$ signs, of which exactly $82-k$ are negative, so
\[
 P_k=(-1)^k.
\]
Cancellation between consecutive products gives
\[
 (x_ix_{82-i})(y_iy_{80-i})=-1
 \qquad(0\le i\le80).
\]
Put $r_i=x_ix_{82-i}$ and $t_i=y_iy_{80-i}$.  Reflection and comparison at
$i$ and $80-i$ give $r_i=r_{i+2}$.  The central products are
$r_{41}=t_{40}=1$, while $r_{40}t_{40}=-1$ gives $r_{40}=-1$.
These two central values fix the odd and even classes of $r_i$; the relation
$r_it_i=-1$ then fixes $t_i$.  This gives \eqref{eq:xreverse} and
\eqref{eq:yreverse}.  In particular $x_0x_{82}=-1$, so
equation~\eqref{eq:fixed-end} gives $uv=1$ and hence $u=v$.
\end{proof}

Separate the even and odd coordinates:
\begin{align*}
 E&=(x_0,x_2,\ldots,x_{82}), & |E|&=42,& E^{\rm rev}&=-E,\\
 O&=(x_1,x_3,\ldots,x_{81}), & |O|&=41,& O^{\rm rev}&=O,\\
 P&=(y_0,y_2,\ldots,y_{80}), & |P|&=41,& P^{\rm rev}&=P,\\
 Q&=(y_1,y_3,\ldots,y_{79}), & |Q|&=40,& Q^{\rm rev}&=-Q,
\end{align*}
and set $B=(u,Q,v)$.  Since $Q$ is skew and $u=v$, the endpoint
contributions to $N_B(k)$ cancel for $1\le k\le40$.  The even-lag equations
in \eqref{eq:fixed-main}, together with \eqref{eq:fixed-end}, become
\[
 N_E(k)+N_B(k)+N_O(k)+N_P(k)=0
 \qquad(1\le k\le41).
\]
Therefore
\[
 (E;B;O;P)\in\BS(42,41).
\]
Its first long sequence is skew and its first short sequence is symmetric,
which is precisely a Turyn sequence in $\TU(41)$ under the standard
definition~\cite{djokovic-quadruples}.  Edmondson, Seberry, and Anderson
exhaustively enumerated all inequivalent Turyn sequences whose long length
is below $43$ and found none of long length $42$~\cite{edmondson}.
This proves Theorem~\ref{thm:fixedq}.  As an independent modern
reproduction of that endpoint, the accompanying outside-in enumerator
exhausts all $461$ normalized depth-five shards for $\TU(41)$:
$57{,}543{,}021$ deterministic nodes and zero solutions.  A separate Python
program exhausts the $2^{19}$ assignments defining the shard cover.

\paragraph{A necessary correction.}
The argument does not follow from $83$ failing to be a sum of two squares.
For $\BS(42,41)$ the relevant row-sum identity at odd parameter $41$ is
\[
 C^2+D^2=4\cdot41-2=162=9^2+9^2.
\]
The final nonexistence step is the Turyn classification, not a two-squares
obstruction.

\section{The raw seed neighbourhood}

Let
\[
 \mathcal E=(A_0;B_0;C_0;D_0)
\]
be the labelled base quadruple obtained from the published $(s,q)$ through
Proposition~\ref{prop:bs}.  We now allow all $334$ signs to change.

\subsection{Margins and endpoint quads}

For a sign sequence $X$, let
\[
 S(X)=\sum_i x_i,\qquad T(X)=\sum_i(-1)^ix_i.
\]
Evaluation of the base-sequence norm identity at $z=1$ gives
\[
 S(A)^2+S(B)^2+S(C)^2+S(D)^2=334,
\]
and evaluation at $z=-1$ gives the same identity for $T$.
The long sums are even, the short sums odd.  Independent sequence
negations and equal-length swaps reduce the ordinary sums to twelve
canonical profiles:
\[
\begin{array}{rrrr}
(4,2,17,5)&(6,0,17,3)&(8,6,15,3)&(10,0,15,3)\\
(10,4,13,7)&(10,8,11,7)&(10,8,13,1)&(12,10,9,3)\\
(14,4,11,1)&(14,8,7,5)&(16,2,7,5)&(18,0,3,1).
\end{array}
\]
Coordinate parity leaves exactly twenty-four compatible alternating
profiles for each ordinary profile, hence $288$ nominal margin shards.
The verifier expands every raw labelled image under the normalizing
negations, swaps, and safe reversals before measuring distance from
$\mathcal E$.

For $1\le k\le82$, the lag-$k$ equation contains
$2(167-2k)$ sign terms, while the lag-$83$ equation contains two.  A zero
sum forces exactly half to be negative; that half-count is odd, so the term
product is $-1$ at every positive lag.  The lag-zero product is $+1$.
Dividing consecutive products telescopes to the standard base-sequence
endpoint parity system, equivalently the endpoint-quad conditions
\begin{align*}
 A_jA_{83-j}B_jB_{83-j}&=
 \begin{cases}-1,&j=0,\\+1,&1\le j\le41,\end{cases}\\
 C_jC_{82-j}D_jD_{82-j}&=+1
 \qquad(0\le j\le40).
\end{align*}
The two central short signs $C_{41}$ and $D_{41}$ are unpaired and remain
unrestricted by these conditions.
The modular seed already has these products.  Thus a repair must flip an
even number of signs in each four-cell endpoint quad.

For one sequence, $(S,T)$ fixes the even- and odd-coordinate sums:
\[
 E_X=\frac{S+T}{2},\qquad O_X=\frac{S-T}{2}.
\]
The minimum number of flips required for a target pair is consequently
exact.  A dynamic program processes the endpoint quads independently,
retaining the cheapest realization of every eight-component margin delta.

\begin{computationallemma}[Dependency-free inner ball]\label{lem:r13}
No exact base sequence lies within raw distance $13$ of $\mathcal E$.
At radius $13$ there are $85$ raw margin targets and none survives the
endpoint-quad dynamic program.  At exact distance $14$, eighteen targets
survive.
\end{computationallemma}

The unique closest raw margin target occurs at distance $8$:
\[
 A:(-18,18),\quad B:(0,0),\quad C:(3,1),\quad D:(-1,-3).
\]
It requires eight positive odd coordinates of $A$ to flip.  They lie in
eight distinct endpoint quads, so the correct quad products cannot be
preserved at distance $8$.  The dynamic program generalizes this observation
through radius $13$ and is independently checked against complete brute
force on small even- and odd-length fixtures.

\subsection{Primitive-root filters}

Write $X(z)=\sum_i x_i z^i$.  For $m\in\{3,4,6\}$, evaluation at a primitive
$m$th root expresses the norm
of each sequence by two integer linear forms $a_m(X),b_m(X)$:
\[
 |X(\zeta_m)|^2
 =a_m(X)^2+\epsilon_m a_m(X)b_m(X)+b_m(X)^2,
 \qquad
 (\epsilon_3,\epsilon_4,\epsilon_6)=(-1,0,1).
\]
Every base sequence satisfies
\begin{equation}
 \sum_{X\in\{A,B,C,D\}}|X(\zeta_m)|^2=334
 \qquad(m=3,4,6).
 \label{eq:smallroots}
\end{equation}
The finite models encode the two linear forms directly and impose their
quadratic norm through a complete allowed-value table.  Terms above $334$
can be discarded because all four norms are nonnegative.

The remaining frontier is split by raw margin target.  Each model fixes all
eight ordinary and alternating class sums, the exact distance interval, every
endpoint-quad product, and \eqref{eq:smallroots}.  For the hard distance-$18$
root cases, endpoint quads are grouped into exact orbits by their residues
modulo $12$ and their oriented seed signs.  Any quotient witness is decoded
to all $334$ physical signs and rechecked with integer arithmetic.

The completed counts are:
\[
\begin{array}{lrrl}
\toprule
\text{scope}&\text{targets}&\text{root-infeasible}&\text{root survivors}\\
\midrule
\text{complete ball through radius }16&197&197&0\\
\text{exact distance-17 shell}&276&276&0\\
\text{exact distance-18 shell}&823&811&12\\
\bottomrule
\end{array}
\]
Solver models with ordinary length-$7$ compression report nine of the twelve
distance-$18$ root-survivor targets infeasible.  Applying the same compression
after coordinate alternation, equivalently primitive-$14$ compression, reports
the final three infeasible.  This supports Claim~\ref{thm:radius} subject to
the proof-certification boundary made explicit next.

\section{Certificates and trust boundary}\label{sec:certification}

The reproducibility package separates four levels of evidence.
\begin{enumerate}
\item Standard-library programs reconstruct all raw margin targets and prove
the radius-$13$ margin-plus-quad exclusion by dynamic programming.
\item Independent artifact checkers reconstruct the $197$, $276$, and $823$
frontiers, verify target uniqueness and completeness, pin all file hashes,
and validate every artifact-selection edge.
\item All twelve decoded root witnesses are independently checked against the
margin, distance, endpoint-quad, and primitive-$3/4/6$ identities.
\item The $1{,}284$ root-infeasible leaves ($197+276+811$), together with the
twelve primitive-$7/14$ compression-infeasible leaves, were originally
reported by one-worker OR-Tools CP-SAT models.  Their JSON records are
reproducible solver outcomes, not independently replayable UNSAT proofs.
\end{enumerate}

Four representative leaves---one each from radius $16$, shell $17$, shell
$18$, and shell-$18$ primitive-$7$ compression---have now been translated
independently to deterministic CNF.  Their pinned binary DRAT traces
regenerate and pass \texttt{drat-trim}; replay of all four took $34.62$
seconds and peaked at $250$ MB resident memory.  A known feasible shell-$18$
root instance also passes an independent positive-model checker with the
symmetry quotient enabled.  All twelve stored root witnesses separately
extend to SAT models of their exactly pinned, unquotiented v2 CNFs and pass
every clause plus independent arithmetic checks.  An exhaustive
contribution-signature regression independently verifies the global root
orbit partition for every even endpoint-quad mask.  This is $4/1{,}296$
proof coverage, not certification of the complete claim.

The release gate for a theorem-level preprint is to replace all $1{,}296$
infeasible solver outcomes in item 4 by
proof-producing SAT or pseudo-Boolean certificates, or by a structurally
independent exact enumeration of every leaf.  Until that gate is met, the
radius-$18$ statement should be read as a strongly audited computational
claim rather than a proof-assistant-level theorem.  The fixed-$q$ reduction
is independent of these solver records.

All production runs were single-threaded.  The original frontier solves
peaked at $176$ MB resident memory; the dependency-free shell-$18$ artifact
checker peaked at $39$ MB.  The independent $\TU(41)$ enumeration completed
$461/461$ shards in $362.92$ seconds of summed serial time; its slowest shard
used $1.38$ MB peak resident memory.  Commands, solver versions, wall times,
memory measurements, and SHA-256 digests are preserved in the repository.

\section{Scope and relation to other work}

The conclusions are deliberately local.
\begin{itemize}
\item Theorem~\ref{thm:fixedq} fixes the exact published $q$.  It does not
exclude changing both $s$ and $q$, a nonspecial Golay quadruple, a Legendre
pair, or a general Hadamard matrix.
\item Claim~\ref{thm:radius} concerns one raw labelled ball.  It does not
exclude the rest of $\BS(84,83)$.
\item A Legendre pair of length $333$ would independently construct an
order-$668$ Hadamard matrix.  Current compression-based work on that route
is mathematically distinct from the seed-repair problem here
~\cite{kotsireas-pq2,chojecki}.
\end{itemize}

Đoković classified $\BS(n+1,n)$ through $n=30$~\cite{djokovic-bs}; that
enumeration does not reach the present $n=83$.  At the audit cutoff, no
public source located by the authors stated either the fixed-$q$ reduction
to $\TU(41)$ or the radius-$18$ repair exclusion.  That is provisional
search evidence, not a guarantee of worldwide priority.  Project policy
prohibits external outreach, so no unpublished-overlap claim is made.

\section*{Code and data availability}

The development repository is
\url{https://github.com/AlecKriebel/Math/tree/main/hadamard_668_search}.
It contains the published seed reconstruction, finite-model generators,
pinned solver records, and independent artifact checkers.  The present
manuscript and certificate work are not yet an immutable public release.
The current reproduction environment uses Python 3.12.13 with OR-Tools
pinned at 9.14.6206; this PDF is built with Tectonic 0.16.9.
Before public circulation they will be frozen at a tagged commit with a
complete manifest and, if size requires, separately archived proof files
identified by SHA-256 digests.

\section*{AI assistance and responsibility}

The derivations, programs, certificates, literature audit, and manuscript
were developed with heavy assistance from ChatGPT 5.6 Sol under Alec
Kriebel's direction.  Alec Kriebel is the named human author and has
requested prominent disclosure that he is not qualified to validate the
mathematics independently.  Independent expert scrutiny is essential.

\begin{thebibliography}{99}
\small
\setlength{\itemsep}{0.25em}

\bibitem{cati-pasechnik}
M. Cati and D. V. Pasechnik,
A database of constructions of Hadamard matrices,
\url{https://arxiv.org/abs/2411.18897}.

\bibitem{chojecki}
P. Chojecki,
Computational search for a Hadamard matrix of order 668 via Legendre pairs
of length 333: status report and roadmap,
March 2026;
\url{https://www.ulam.ai/research/frontier-had.pdf}.

\bibitem{djokovic-quadruples}
D. Ž. Đoković,
Aperiodic complementary quadruples of binary sequences,
\emph{J. Combin. Math. Combin. Comput.} 27 (1998), 3--31;
\url{https://combinatorialpress.com/article/jcmcc/Volume%20027/vol-27-paper%201.pdf}.

\bibitem{djokovic-bs}
D. Ž. Đoković,
Classification of base sequences $\BS(n+1,n)$,
\emph{Int. J. Combin.} (2010), Article ID 851857;
\url{https://arxiv.org/abs/1002.1414}.

\bibitem{edmondson}
G. M. Edmondson, J. Seberry, and M. R. Anderson,
On the existence of Turyn sequences of length less than 43,
\emph{Math. Comp.} 62 (1994), 351--362;
\url{https://doi.org/10.1090/S0025-5718-1994-1203733-8}.

\bibitem{eliahou-2025}
S. Eliahou,
A 64-modular Hadamard matrix of order 668,
\emph{Australas. J. Combin.} 93(2) (2025), 422--427;
\url{https://ajc.maths.uq.edu.au/pdf/93/ajc_v93_p422.pdf}.

\bibitem{eliahou-2026}
S. Eliahou,
An update on modular Hadamard matrices,
\emph{J. Algebraic Combin.} 64 (2026), Article 1;
\url{https://doi.org/10.1007/s10801-026-01544-5}.

\bibitem{kotsireas-pq2}
I. Kotsireas, R. Gallardo-Cava, A. I. Gómez, and D. Gómez-Pérez,
On the search of binary Legendre pairs of length $pq^2$,
\emph{J. Symbolic Comput.} 138 (2027), Article 102606;
\url{https://doi.org/10.1016/j.jsc.2026.102606}.

\bibitem{sage}
SageMath,
Hadamard matrix constructions and current unknown orders;
\url{https://doc.sagemath.org/html/en/reference/combinat/sage/combinat/matrices/hadamard_matrix.html}.

\end{thebibliography}

\end{document}
